\textbf{Phần 1. Trắc nghiệm lựa chọn}
\setcounterpageref{ex}{0}

\begin{ex} %[2H5N28]
Hình nào sau đây có 4 đỉnh và 4 cạnh?
\choice
{Hình tam giác}
{Hình tròn}
{Đường thẳng}
{\True Hình tứ giác}
\loigiai{
Hình tứ giác có 4 đỉnh và 4 cạnh.\\
\textbf{Đáp án của câu hỏi là D}
}
\end{ex}

\begin{ex} %[2H5N28]
Đoạn thẳng $PQ$ dài $15\text{ cm}$, đoạn thẳng $RS$ dài $10\text{ cm}$. Cả hai đoạn thẳng dài là
\choice
{$20\text{ cm}$}
{\True $25\text{ cm}$}
{$30\text{ cm}$}
{$35\text{ cm}$}
\loigiai{
$15 + 10 = 25\text{ cm}$.\\
\textbf{Đáp án của câu hỏi là B}
}
\end{ex}

\begin{ex} %[2H5H28]
Đường gấp khúc $ABCD$ gồm các đoạn thẳng $AB = 5\text{ cm}, BC = 6\text{ cm}, CD = 7\text{ cm}$. Độ dài đường gấp khúc là
\choice
{\True $18\text{ cm}$}
{$17\text{ cm}$}
{$19\text{ cm}$}
{$20\text{ cm}$}
\loigiai{
$5 + 6 + 7 = 18\text{ cm}$.\\
\textbf{Đáp án của câu hỏi là A}
}
\end{ex}

\begin{ex} %[2H5H28]
Hình tứ giác có độ dài các cạnh lần lượt là $4\text{ cm}, 5\text{ cm}, 6\text{ cm}, 7\text{ cm}$. Tổng độ dài các cạnh là
\choice
{$21\text{ cm}$}
{\True $22\text{ cm}$}
{$23\text{ cm}$}
{$24\text{ cm}$}
\loigiai{
$4 + 5 + 6 + 7 = 22\text{ cm}$.\\
\textbf{Đáp án của câu hỏi là B}
}
\end{ex}

\begin{ex} %[2H5H28]
Cho 3 điểm $A, B, C$ thẳng hàng theo thứ tự đó. Biết $AB = 7\text{ cm}, BC = 8\text{ cm}$. Độ dài $AC$ là
\choice
{\True $15\text{ cm}$}
{$14\text{ cm}$}
{$16\text{ cm}$}
{$17\text{ cm}$}
\loigiai{
$AC = 7 + 8 = 15\text{ cm}$.\\
\textbf{Đáp án của câu hỏi là A}
}
\end{ex}

\begin{ex} %[2H5H28]
Đoạn thẳng $CD$ dài $20\text{ cm}$, điểm $I$ là điểm chính giữa của đoạn thẳng $CD$. Độ dài đoạn $CI$ là
\choice
{$8\text{ cm}$}
{$12\text{ cm}$}
{\True $10\text{ cm}$}
{$15\text{ cm}$}
\loigiai{
$20 : 2 = 10\text{ cm}$.\\
\textbf{Đáp án của câu hỏi là C}
}
\end{ex}

\begin{ex} %[2H5V28]
Đường gấp khúc gồm 4 đoạn thẳng có độ dài bằng nhau, mỗi đoạn dài $8\text{ cm}$. Độ dài đường gấp khúc đó là
\choice
{$30\text{ cm}$}
{$34\text{ cm}$}
{$36\text{ cm}$}
{\True $32\text{ cm}$}
\loigiai{
$8 \times 4 = 32\text{ cm}$.\\
\textbf{Đáp án của câu hỏi là D}
}
\end{ex}

\begin{ex} %[2H5V28]
Một hình tứ giác có 4 cạnh dài bằng nhau, tổng độ dài 4 cạnh là $28\text{ cm}$. Độ dài mỗi cạnh của tứ giác là
\choice
{$6\text{ cm}$}
{$8\text{ cm}$}
{\True $7\text{ cm}$}
{$9\text{ cm}$}
\loigiai{
$28 : 4 = 7\text{ cm}$.\\
\textbf{Đáp án của câu hỏi là C}
}
\end{ex}

\begin{ex} %[2H5V28]
Có 10 que tính, dùng để xếp được 1 hình vuông và mấy hình tam giác rời nhau?
\choice
{$1$ hình tam giác}
{$3$ hình tam giác}
{$4$ hình tam giác}
{\True $2$ hình tam giác}
\loigiai{
Hình vuông dùng 4 que, còn $10 - 4 = 6$ que tính.\\
$6 : 3 = 2$ hình tam giác.\\
\textbf{Đáp án của câu hỏi là D}
}
\end{ex}

\begin{ex} %[2H5V28]
Đường gấp khúc $ABC$ dài $25\text{ cm}$. Biết đoạn thẳng $AB$ dài hơn đoạn thẳng $BC$ là $5\text{ cm}$. Độ dài đoạn thẳng $BC$ là
\choice
{$8\text{ cm}$}
{\True $10\text{ cm}$}
{$12\text{ cm}$}
{$15\text{ cm}$}
\loigiai{
Đoạn $BC$: $(25 - 5) : 2 = 10\text{ cm}$.\\
\textbf{Đáp án của câu hỏi là B}
}
\end{ex}

\begin{ex} %[2H5V28]
Cho đoạn thẳng $AB = 18\text{ cm}$, vẽ thêm đoạn thẳng $BC = 7\text{ cm}$ thẳng hàng với $AB$. Đoạn thẳng $AC$ dài là
\choice
{\True $25\text{ cm}$}
{$24\text{ cm}$}
{$26\text{ cm}$}
{$27\text{ cm}$}
\loigiai{
$18 + 7 = 25\text{ cm}$.\\
\textbf{Đáp án của câu hỏi là A}
}
\end{ex}

\begin{ex} %[2H5C28]
Một mảnh giấy hình vuông được cắt thành 4 hình tam giác bằng nhau. Ghép 4 hình tam giác đó lại có thể tạo thành hình nào?
\choice
{Chỉ tạo thành hình tròn}
{Không thể ghép lại}
{\True Hình vuông lớn hoặc hình chữ nhật hoặc hình tam giác lớn}
{Chỉ tạo thành đường thẳng}
\loigiai{
Có thể tạo thành hình vuông, chữ nhật hoặc tam giác lớn.\\
\textbf{Đáp án của câu hỏi là C}
}
\end{ex}

\textbf{Phần 2. Trắc nghiệm đúng sai}
\setcounterpageref{ex}{0}

\begin{ex} %[2H5H28]
Xét tính đúng sai của các nhận định hình học:
\choiceTF
{\True Hình tam giác có 3 cạnh và 3 đỉnh.}
{\True Hình tứ giác có 4 cạnh và 4 đỉnh.}
{Đoạn thẳng không giới hạn ở hai đầu.}
{\True Đoạn thẳng bị giới hạn bởi 2 điểm đầu mút.}
\loigiai{
\begin{itemize}
    \item Tam giác: 3 cạnh, 3 đỉnh: Đúng.
    \item Tứ giác: 4 cạnh, 4 đỉnh: Đúng.
    \item Đoạn thẳng bị giới hạn 2 đầu. Ý 3 Sai, Ý 4 Đúng.
\end{itemize}
\begin{itemize}
    \item a) \textbf{Đúng}.
    \item b) \textbf{Đúng}.
    \item c) \textbf{Sai}.
    \item d) \textbf{Đúng}.
\end{itemize}
}
\end{ex}

\begin{ex} %[2H5V28]
Đường gấp khúc $MNPQ$ có $MN = 8\text{ cm}, NP = 10\text{ cm}, PQ = 12\text{ cm}$. Xét tính đúng sai:
\choiceTF
{\True Độ dài đường gấp khúc $MNPQ$ là $30\text{ cm}$ ($8 + 10 + 12 = 30$).}
{\True Đoạn $PQ$ dài hơn đoạn $MN$ là $4\text{ cm}$ ($12 - 8 = 4$).}
{Đường gấp khúc $MNPQ$ gồm 4 đoạn thẳng.}
{\True Đường gấp khúc $MNPQ$ gồm 3 đoạn thẳng là $MN, NP, PQ$.}
\loigiai{
\begin{itemize}
    \item Độ dài: $8 + 10 + 12 = 30\text{ cm}$: Đúng.
    \item $PQ - MN = 12 - 8 = 4\text{ cm}$: Đúng.
    \item Gồm 3 đoạn thẳng (không phải 4). Ý 3 Sai, Ý 4 Đúng.
\end{itemize}
\begin{itemize}
    \item a) \textbf{Đúng}.
    \item b) \textbf{Đúng}.
    \item c) \textbf{Sai}.
    \item d) \textbf{Đúng}.
\end{itemize}
}
\end{ex}

\textbf{Phần 3. Trả lời ngắn}
\setcounterpageref{ex}{0}

\begin{ex} %[2H5H28]
Tính độ dài đường gấp khúc gồm 3 đoạn thẳng lần lượt dài $6\text{ cm}, 8\text{ cm}, 10\text{ cm}$ theo xăng-ti-mét.
\shortans{24}
\loigiai{
\textbf{Đáp số:} $24$\\
$6 + 8 + 10 = 24\text{ cm}$.
}
\end{ex}

\begin{ex} %[2H5H28]
Một tứ giác có độ dài 4 cạnh là $5\text{ cm}, 7\text{ cm}, 9\text{ cm}, 11\text{ cm}$. Tính tổng độ dài các cạnh theo xăng-ti-mét.
\shortans{32}
\loigiai{
\textbf{Đáp số:} $32$\\
$5 + 7 + 9 + 11 = 32\text{ cm}$.
}
\end{ex}

\begin{ex} %[2H5V28]
Đường gấp khúc gồm 5 đoạn thẳng bằng nhau, mỗi đoạn dài $4\text{ cm}$. Tính độ dài đường gấp khúc theo xăng-ti-mét.
\shortans{20}
\loigiai{
\textbf{Đáp số:} $20$\\
$4 \times 5 = 20\text{ cm}$.
}
\end{ex}

\begin{ex} %[2H5V28]
Cho 3 điểm thẳng hàng $A, B, C$ với $AB = 12\text{ cm}, AC = 20\text{ cm}$ ($B$ nằm giữa $A$ và $C$). Tính độ dài $BC$ theo xăng-ti-mét.
\shortans{8}
\loigiai{
\textbf{Đáp số:} $8$\\
$BC = 20 - 12 = 8\text{ cm}$.
}
\end{ex}

\begin{ex} %[2H5V28]
Có 15 que tính, có thể xếp được nhiều nhất bao nhiêu hình tam giác rời nhau?
\shortans{5}
\loigiai{
\textbf{Đáp số:} $5$\\
$15 : 3 = 5$ hình.
}
\end{ex}

\begin{ex} %[2H5C28]
Một hình tứ giác có tổng độ dài 4 cạnh là $36\text{ cm}$, biết 4 cạnh có độ dài bằng nhau. Tính độ dài mỗi cạnh theo xăng-ti-mét.
\shortans{9}
\loigiai{
\textbf{Đáp số:} $9$\\
$36 : 4 = 9\text{ cm}$.
}
\end{ex}
